\documentclass[10pt,a4paper]{amsart}
\usepackage{amsmath,amssymb,amsthm,graphicx}
\usepackage{verbatim}
\newenvironment{reference}{\comment}{\endcomment}
%\newenvironment{reference}{}{}
\newenvironment{slogan}{\comment}{\endcomment}
\newenvironment{history}{\comment}{\endcomment}

% For commutative diagrams we use Xy-pic
\usepackage[all]{xy}

% We use 2cell for 2-commutative diagrams.
\xyoption{2cell}
% Theorem environments.
%
\theoremstyle{plain}
\newtheorem{theorem}[subsection]{Theorem}
\newtheorem{proposition}[subsection]{Proposition}
\newtheorem{lemma}[subsection]{Lemma}

\theoremstyle{definition}
\newtheorem{definition}[subsection]{Definition}
\newtheorem{example}[subsection]{Example}
\newtheorem{exercise}[subsection]{Exercise}
\newtheorem{situation}[subsection]{Situation}

\theoremstyle{remark}
\newtheorem{remark}[subsection]{Remark}
\newtheorem{remarks}[subsection]{Remarks}

\numberwithin{equation}{subsection}

% Macros
%
\def\lim{\mathop{\mathrm{lim}}\nolimits}
\def\colim{\mathop{\mathrm{colim}}\nolimits}
\def\Spec{\mathop{\mathrm{Spec}}}
\def\Hom{\mathop{\mathrm{Hom}}\nolimits}
\def\Ext{\mathop{\mathrm{Ext}}\nolimits}
\def\SheafHom{\mathop{\mathcal{H}\!\mathit{om}}\nolimits}
\def\SheafExt{\mathop{\mathcal{E}\!\mathit{xt}}\nolimits}
\def\Sch{\mathit{Sch}}
\def\Mor{\mathop{\mathrm{Mor}}\nolimits}
\def\Ob{\mathop{\mathrm{Ob}}\nolimits}
\def\Sh{\mathop{\mathit{Sh}}\nolimits}
\def\NL{\mathop{N\!L}\nolimits}
\def\CH{\mathop{\mathrm{CH}}\nolimits}
\def\proetale{{pro\text{-}\acute{e}tale}}
\def\etale{{\acute{e}tale}}
\def\QCoh{\mathit{QCoh}}
\def\Ker{\mathop{\mathrm{Ker}}}
\def\Im{\mathop{\mathrm{Im}}}
\def\Coker{\mathop{\mathrm{Coker}}}
\def\Coim{\mathop{\mathrm{Coim}}}

% Boxtimes
%
\DeclareMathSymbol{\boxtimes}{\mathbin}{AMSa}{"02}

%
% Macros for moduli stacks/spaces
%
\def\QCohstack{\mathcal{QC}\!\mathit{oh}}
\def\Cohstack{\mathcal{C}\!\mathit{oh}}
\def\Spacesstack{\mathcal{S}\!\mathit{paces}}
\def\Quotfunctor{\mathrm{Quot}}
\def\Hilbfunctor{\mathrm{Hilb}}
\def\Curvesstack{\mathcal{C}\!\mathit{urves}}
\def\Polarizedstack{\mathcal{P}\!\mathit{olarized}}
\def\Complexesstack{\mathcal{C}\!\mathit{omplexes}}
% \Pic is the operator that assigns to X its picard group, usage \Pic(X)
% \Picardstack_{X/B} denotes the Picard stack of X over B
% \Picardfunctor_{X/B} denotes the Picard functor of X over B
\def\Pic{\mathop{\mathrm{Pic}}\nolimits}
\def\Picardstack{\mathcal{P}\!\mathit{ic}}
\def\Picardfunctor{\mathrm{Pic}}
\def\Deformationcategory{\mathcal{D}\!\mathit{ef}}
\usepackage{cleveref}
\crefname{originallemma}{original lemma}{original lemmas}
\Crefname{originallemma}{Original lemma}{Original lemmas}
\makeatletter
\newlabel{algebra-lemma-quasi-compact}{{00E8}{source}{}{sourcepointer-1}{}}
\newlabel{algebra-lemma-quasi-compact@cref}{{[originallemma][0][]00E8}{[1][0][]source}{}{}{}}
\makeatother
\crefname{originaldefinition}{original definition}{original definitions}
\Crefname{originaldefinition}{Original definition}{Original definitions}
\makeatletter
\newlabel{sites-definition-quasi-compact}{{090H}{source}{}{sourcepointer-2}{}}
\newlabel{sites-definition-quasi-compact@cref}{{[originaldefinition][0][]090H}{[1][0][]source}{}{}{}}
\makeatother
\title{An infinite flat descent family with the equalizer property but no finite subfamily}
\author{The Stacks Project Authors}
\date{stacks-project:examples.tex:local-snapshot-20261012; GNU FDL 1.2 or later}
\begin{document}
\maketitle
\noindent\textit{Editorial note.} All delivered mathematical text is the author's own, copied from the pinned Stacks Project chapter file examples.tex; no statement or proof step was rewritten, reordered or removed.\par\medskip
\noindent\textit{Editorial note.} The counted proof is the author's complete argument: Claim I of the canonical-topology section; the companion lemma lemma-Z-not-quasi-compact consumes this family together with Claim II.\par\medskip
\noindent\textit{Editorial note.} The label lemma-non-fpqc-descent names the counted statement, and the proof environment immediately following it in the source is delivered verbatim as its proof.\par\medskip
\noindent\textit{Editorial note.} The counted statement is carried by the author's own lemma environment in the chapter source.\par\medskip
\noindent\textit{Allowed standard prerequisites (editorial).}\begin{itemize}
\item A nonprincipal ultrafilter U on the set P of prime numbers: every subset of P or its complement lies in U, U is closed under finite intersection, and no finite set lies in U
\item For A in U the ring Z-A is the localization of the integers at the multiplicative set generated by the primes in A; its spectrum is the generic point plus the closed primes outside A
\item Localization is exact and commutes with tensor products; a torsion abelian group decomposes as the direct sum of its p-primary components
\item The equalizer condition for a family of flat ring maps A to A-i on modules: M equals the equalizer of the two natural maps from the product of M tensor A-i to the product over pairs of M tensor A-i tensor A A-j
\end{itemize}\medskip
\ifdefined\hypertarget\hypertarget{sourcepointer-1}{}\fi
\noindent\textit{Source reference.} Original lemma 00E8: sources/stacks-chow/tags.txt, source lines 525--525. Prerequisite pointer invoked at 7142; it names a standard fact of a sibling chapter that the counted argument uses, and every step of the counted proof built on it is retained verbatim in the delivered body.\par\medskip
\ifdefined\hypertarget\hypertarget{sourcepointer-2}{}\fi
\noindent\textit{Source reference.} Original definition 090H: sources/stacks-chow/tags.txt, source lines 11410--11410. Prerequisite pointer invoked at 7146; it names a standard fact of a sibling chapter that the counted argument uses, and every step of the counted proof built on it is retained verbatim in the delivered body.\par\medskip
the spectrum of the integers as a topological space,
because any spectrum is quasi-compact as a topological space
(Algebra, Lemma \ref{algebra-lemma-quasi-compact}).
No, it refers to the spectrum of the integers in the
canonical topology on the category of schemes, and the
definition of a quasi-compact object in a site
(Sites, Definition \ref{sites-definition-quasi-compact}).

\medskip\noindent
Let $U$ be a nonprincipal ultrafilter on the set $P$ of prime numbers.
For a subset $T \subset P$ we denote $T^c = P \setminus T$ the complement.
For $A \in U$ let $S_A \subset \mathbf{Z}$ be the multiplicative
subset generated by $p \in A$. Set
$$
\mathbf{Z}_A = S_A^{-1}\mathbf{Z}
$$
Observe that $\Spec(\mathbf{Z}_A) = \{(0)\} \cup A^c \subset \Spec(\mathbf{Z})$
if we think of $P$ as the set of closed points of $\Spec(\mathbf{Z})$.
If $A, B \in U$, then $A \cap B \in U$ and $A \cup B \in U$
and we have
$$
\mathbf{Z}_{A \cap B} =
\mathbf{Z}_A \times_{\mathbf{Z}_{A \cup B}} \mathbf{Z}_B
$$
(fibre product of rings). In particular, for any
integer $n$ and elements $A_1, \ldots, A_n \in U$ the morphisms
$$
\Spec(\mathbf{Z}_{A_1}) \amalg \ldots \amalg \Spec(\mathbf{Z}_{A_n})
\longrightarrow \Spec(\mathbf{Z})
$$
factors through $\Spec(\mathbf{Z}[1/p])$ for some $p$
(namely for any $p \in A_1 \cap \ldots \cap A_n$).
We conclude that the family of flat morphisms
$\{\Spec(\mathbf{Z}_A) \to \Spec(\mathbf{Z})\}_{A \in U}$ is jointly
surjective, but no finite subset is.

\medskip\noindent
For a $\mathbf{Z}$-module $M$ we set
$$
M_A = S_A^{-1}M = M \otimes_{\mathbf{Z}} \mathbf{Z}_A
$$
Claim I: for every $\mathbf{Z}$-module $M$ we have
$$
M =
\text{Equalizer}\left(
\xymatrix{
\prod\nolimits_{A \in U} M_A \ar@<1ex>[r] \ar@<-1ex>[r] &
\prod\nolimits_{A, B \in U} M_{A \cup B}
}
\right)
$$
First, assume $M$ is torsion free. Then $M_A \subset M_P$ for all $A \in U$.
Hence we see that we have to prove
$$
M = \bigcap\nolimits_{A \in U} M_A\text{ inside }M_P = M \otimes \mathbf{Q}
$$
Namely, since $U$ is nonprincipal, for any prime $p$ we have
$\{p\}^c \in U$. Also, $M_{\{p\}^c} = M_{(p)}$ is equal to the
localization at the prime $(p)$.
Thus the above is clear because already $M_{(2)} \cap M_{(3)} = M$.
Next, assume $M$ is torsion. Then we have
$$
M = \bigoplus\nolimits_{p \in P} M[p^\infty]
$$
and correspondingly we have
$$
M_A = \bigoplus\nolimits_{p \not \in A} M[p^\infty]
$$
because we are localizing at the primes in $A$. Suppose that
$(x_A) \in \prod M_A$ is in the equalizer. Denote
$x_p = x_{\{p\}^c} \in M[p^{\infty}]$. Then the equalizer
property says
$$
x_A = (x_p)_{p \not \in A}
$$
and in particular it says that $x_p$ is zero for all but a finite
number of $p \not \in A$. To finish the proof in the torsion case
it suffices to show that $x_p$ is zero for all but a finite number
of primes $p$. If not write $\{p \in P \mid x_p \not = 0\} = T \amalg T'$
as the disjoint union of two infinite sets. Then either
$T \not \in U$ or $T' \not \in U$ because $U$ is an ultrafilter
(namely if both $T, T'$ are in $U$ then $U$ contains $T \cap T' = \emptyset$
which is not allowed). Say $T \not \in U$. Then $T = A^c$ and this
contradicts the finiteness mentioned above.
Finally, suppose that $M$ is a general module. Then we look
at the short exact sequence
$$
0 \to M_{tors} \to M \to M/M_{tors} \to 0
$$
and we look at the following large diagram
$$
\xymatrix{
M_{tors} \ar[r] \ar[d] &
\prod\nolimits_{A \in U} M_{tors, A} \ar@<1ex>[r] \ar@<-1ex>[r] \ar[d] &
\prod\nolimits_{A, B \in U} M_{tors, A \cup B} \ar[d] \\
M \ar[r] \ar[d] &
\prod\nolimits_{A \in U} M_A \ar@<1ex>[r] \ar@<-1ex>[r] \ar[d] &
\prod\nolimits_{A, B \in U} M_{A \cup B} \ar[d] \\
M/M_{tors} \ar[r] &
\prod\nolimits_{A \in U} (M/M_{tors})_A \ar@<1ex>[r] \ar@<-1ex>[r] &
\prod\nolimits_{A, B \in U} (M/M_{tors})_{A \cup B} \\
}
$$
Doing a diagram chase using exactness of the columns and the result
for the torsion module $M_{tors}$ and the torsion free module
$M/M_{tors}$ proving Claim I for $M$.

\begin{lemma}
\label{lemma-non-fpqc-descent}
There exists a ring $A$ and an infinite family of flat ring maps
$\{A \to A_i\}_{i \in I}$ such that for every $A$-module $M$ 
$$
M =
\text{Equalizer}\left(
\xymatrix{
\prod\nolimits_{i \in I} M \otimes_A A_i \ar@<1ex>[r] \ar@<-1ex>[r] &
\prod\nolimits_{i, j \in I} M \otimes_A A_i \otimes_A A_j
}
\right)
$$
but there is no finite subfamily where the same thing is true.
\end{lemma}

\begin{proof}
See discussion above.
\end{proof}

\end{document}
